\begin{abstractpage}

\begin{abstract}{romanian}


Scopul acestei lucrări este de a prezenta proiectarea și implementarea unui \textbf{asistent AI desktop conștient de emoții, cu recunoaștere facială a emoțiilor}. Aplicația îi permite utilizatorului să interacționeze cu asistentul prin text și voce, folosind în același timp recunoașterea facială a emoțiilor pentru a îmbunătăți contextul interacțiunii. În loc să se bazeze doar pe ceea ce utilizatorul scrie sau spune, asistentul ia în considerare și modul în care utilizatorul pare să se simtă în timpul conversației.

Pentru a atinge acest obiectiv, aplicația este construită ca un sistem modular care conectează mai multe componente de inteligență artificială într-un singur asistent desktop. Sistemul combină procesarea intrărilor utilizatorului, analiza emoțiilor faciale, transcrierea vorbirii, generarea răspunsurilor și comunicarea în timp real între module. Starea emoțională detectată este folosită ca un context suplimentar pentru asistent, permițând ca răspunsurile generate să fie mai bine adaptate la interacțiunea curentă.

Soluția este formată din cinci componente principale:

\begin{itemize}
    \item O interfață de asistent desktop care îi permite utilizatorului să interacționeze cu sistemul prin text și voce;

    \item O componentă backend responsabilă de gestionarea sesiunii de interacțiune și a comunicării dintre modulele aplicației;

    \item O componentă speech-to-text utilizată pentru convertirea intrării vocale a utilizatorului în text;

    \item O componentă de recunoaștere facială a emoțiilor utilizată pentru detectarea stării emoționale aparente a utilizatorului din imaginile capturate de webcam;

    \item O componentă de generare a răspunsurilor care primește mesajul utilizatorului împreună cu contextul emoțional detectat și generează răspunsul asistentului.
\end{itemize}

Lucrarea se concentrează pe prezentarea celor cinci componente și a modului în care acestea sunt conectate în aplicația finală, precum și pe antrenarea și evaluarea modelului de recunoaștere facială a emoțiilor, integrarea contextului emoțional în răspunsurile asistentului și prezentarea rezultatelor experimentale obținute.

\end{abstract}

\newpage

\begin{abstract}{english}

The goal of this paper is to present the design and implementation of an \textbf{Emotion-Aware Desktop AI Assistant with Facial Emotion Recognition}. The application allows the user to interact with an assistant through text and voice, while also using facial emotion recognition to improve the context of the interaction. Instead of relying only on what the user writes or says, the assistant also takes into account how the user appears to feel during the conversation.

To achieve this goal, the application is built as a modular system that connects several artificial intelligence components into a single desktop assistant. The system combines user input processing, facial emotion analysis, speech transcription, response generation and real-time communication between the modules. The detected emotional state is used as additional context for the assistant, allowing the generated answers to be more adapted to the current interaction.

The solution consists of five main components:

\begin{itemize}
    \item A desktop assistant interface that allows the user to interact with the system through text and voice;

    \item A backend component responsible for managing the interaction session and the communication between the application modules;

    \item A speech-to-text component used for converting the user's voice input into text;

    \item A facial emotion recognition component used for detecting the user's apparent emotional state from webcam images;

    \item A response generation component that receives the user's message together with the detected emotional context and generates the assistant's answer.
\end{itemize}

The paper focuses on the presentation of the five components and their connection in the final application, as well as the training and evaluation of the facial emotion recognition model, the integration of emotional context into the assistant's responses, and the presentation of the experimental results obtained.

\end{abstract}

\end{abstractpage}
